\scanpage{88}
les fibres de $\unsure{E|\mathcal E}/B$), alors $\exists$ adjoint à droite $F_+^c$ de $F_c^*$.

\textcircled{3} \uline{Prop}. Supposons que pour \unsure{tt|tout} $b\in\mathrm{Ob}\,B$, la limite
inductive $\underset{c\in\mathrm{Ob}\,{}^b/B}{\underline{\mathrm{Lim}}}\ f^*(X(c))$ dans $E_b$ existe,
et supposons que cette $\underline{\lim}$ commute
à tout \unsure{ch.gt|ch.g\textsuperscript{t}} de base $u^*$ ($a:b'\mathbin{\unsure{-|\to}}b$).
Alors le foncteur adjoint à gauche $\deleted{F}i_*^{\deleted{\unsure{c|\varepsilon}}}$
de \deleted{$\unsure{F'|E}$} $i:\underleftarrow{\mathrm{Lim}}\ E/B\longrightarrow\underline{\Gamma}(B/E)$ existe,
et est donné par

\[
\deleted{F}i_*^{\deleted{\unsure{c|\varepsilon}}}(X)(b)
=\underset{x\in\mathrm{Ob}\,{}^b/B}{\underrightarrow{\mathrm{Lim}}}\ f^*(X(x))
\]

[Le morphisme d'adjonction

\[
\deleted{F}i_*^{\deleted{\unsure{c|\varepsilon}}}\,\deleted{F}i_{\deleted{\unsure{*|c}}}^{\deleted{\unsure{*|\varepsilon}}}(X)\longrightarrow X
\]

étant le morphisme évident].

\inserted{Si de plus $E\longrightarrow X$ est cofibrant, on trouve des adjoints à gauche des $b_c^*:\underleftarrow{\mathrm{Lim}}\ E/B\mathbin{\unsure{-|\to}}E_b$}

% slanted
\margin{Énoncé dual pour la construction d'un adjoint \uline{à droite}, pour
$\boxed{(\underleftarrow{\mathrm{Lim}}\ E/B)^\circ\to\underline{\Gamma}(E^\triangle/B)^\circ}$
moyennant la comm. des $f^*$ aux petites $\underleftarrow{\lim}$. Alors les adj.}

\deleted{N.B. Ce critère ne fait pas double emploi avec
l'hyp. \unsure{ess|ass} de commutation des $u^*$
aux $\underrightarrow{\lim}$ n'étant pas dans la $=$
direction que celle que les $u^*$ aient des
adjts à g. $u_*$ -- qui implique \deleted{\unsure{nt}} que $u^*$ commute
aux $\underleftarrow{\lim}$ \ldots\ ].}

\margin{\uline{à droite} des $b_c^*:\underleftarrow{\mathrm{Lim}}\ E/B\longrightarrow E_b$ existent également et se construisent explicit\textsuperscript{t}, s'il en est ainsi des $f_*^{\unsure{s|+|*}}$}

% archive sheet number: 87
